\section{AMC 8 1990}
        \begin{enumerate}
        	\item (\textit{AMC 8 1990}) What is the smallest sum of two $3$-digit numbers that can be obtained by placing each of the six digits $4,5,6,7,8,9$ in one of the six boxes in this addition problem?


 
                    \begin{figure}[H]
                        \centering
                        \includegraphics[width=0.5\textwidth]{Problem/AMC_8_AJHSME_1990_1_Problem_Figure_1.png}
                    \end{figure}
                    
$\text{(A)}\ 947 \qquad \text{(B)}\ 1037 \qquad \text{(C)}\ 1047 \qquad \text{(D)}\ 1056 \qquad \text{(E)}\ 1245$

 \textbf{Answer: }C



\textbf{Solution 1}

The minimum sum is found by choosing smaller digits to have higher place values.  In this case, $4$ and $5$ should be hundreds digits, $6$ and $7$ should be tens digits, and $8$ and $9$ should be units digits.  The minimum sum is therefore $468 + 579 = \boxed{\textbf{(C)}\ 1047}$.

 



\textbf{Solution 2}

Let the two $3$-digit numbers be $\overline{abc}$ and $\overline{def}$.  Their sum is $100(a+d)+10(b+e)+(c+f)$.  

 To minimize this sum, $a+d$ should be minimized, which happens if $\{a, d\}=\{4,5\}$.  By similar logic using the remaining digits, the sum is minimized if $\{b,e\}=\{6,7\}$ while $\{c,f\}=\{8,9\}$.  The minimum sum is therefore $100(9)+10(13)+(17)=\boxed{\textbf{(C)}\ 1047}$



	\item (\textit{AMC 8 1990}) Which digit of $.12345$, when changed to $9$, gives the largest number?


 $\text{(A)}\ 1 \qquad \text{(B)}\ 2 \qquad \text{(C)}\ 3 \qquad \text{(D)}\ 4 \qquad \text{(E)}\ 5$

 \textbf{Answer: }A



\textbf{Solution 1}

Since $\boxed{\textbf{(A)}\ 1}$ has the highest place value, changing it to $9$ will result in the largest number.

 



\textbf{Solution 2}

Changing $1$ to $9$ results in a number greater than $.9$ while changing any other digit to $9$ results in a number less than $.2$.  Therefore, changing $\boxed{\textbf{(A)}\ 1}$ results in the largest number.

 


 


	\item (\textit{AMC 8 1990}) What fraction of the square is shaded?


 
                    \begin{figure}[H]
                        \centering
                        \includegraphics[width=0.5\textwidth]{Problem/AMC_8_AJHSME_1990_3_Problem_Figure_1.png}
                    \end{figure}
                    
$\text{(A)}\ \frac{1}{3} \qquad \text{(B)}\ \frac{2}{5} \qquad \text{(C)}\ \frac{5}{12} \qquad \text{(D)}\ \frac{3}{7} \qquad \text{(E)}\ \frac{1}{2}$

 \textbf{Answer: }E



\textbf{Solution 1}


                    \begin{figure}[H]
                        \centering
                        \includegraphics[width=0.5\textwidth]{Solution/AMC_8_AJHSME_1990_3_Solution_Figure_1.png}
                    \end{figure}
                    

 Swap the two trapezoids in the middle layer, as shown.  The shaded region and unshaded region now form congruent triangles, $\boxed{\textbf{(E)}\ \frac{1}{2}}$ of the square is shaded.

 



\textbf{Solution 2}

Notice that the figure is made of $3$ pairs of congruent shapes, and each of these pairs has $1$ shaded and $1$ unshaded shape.  Therefore, $\boxed{\textbf{(E)}\ \frac{1}{2}}$ of the square is shaded.

 


	\item (\textit{AMC 8 1990}) Which of the following could \textbf{not} be the unit's digit [one's digit] of the square of a whole number?


 $\text{(A)}\ 1 \qquad \text{(B)}\ 4 \qquad \text{(C)}\ 5 \qquad \text{(D)}\ 6 \qquad \text{(E)}\ 8$

 \textbf{Answer: }E



\textbf{Solution 1}

Let $n$ be a whole number.

 If $n$ ends in $0$, then $n^2$ ends in $0$.

 If $n$ ends in $1$ or $9$, then $n^2$ ends in $1$.

 If $n$ ends in $2$ or $8$, then $n^2$ ends in $4$.

 If $n$ ends in $3$ or $7$, then $n^2$ ends in $9$.

 If $n$ ends in $4$ or $6$, then $n^2$ ends in $6$.

 If $n$ ends in $5$, then $n^2$ ends in $5$.

 Therefore, $n^2$ cannot end in $\boxed{\textbf{(E)}\ 8}$.

 



\textbf{Solution 2}

Notice that $1^2=1$, $2^2=4$, $5^2=25$, and $4^2=16$, to eliminate (A), (B), (C), and (D) respectively.  So $\boxed{\textbf{(E)}\ 8}$ must be correct.

 


	\item (\textit{AMC 8 1990}) Which of the following is closest to the product $(.48017)(.48017)(.48017)$?


 $\text{(A)}\ 0.011 \qquad \text{(B)}\ 0.110 \qquad \text{(C)}\ 1.10 \qquad \text{(D)}\ 11.0 \qquad \text{(E)}\ 110$

 \textbf{Answer: }B



\textbf{Solution 1}

Since $.48017$ is between $.4$ and $.5$, $(.48017)^3$ must be between $(.4)^3 = .064$ and $(.5)^3 = .125$.  So $.40817$ must be closest to $\boxed{\textbf{(B)}\ 0.110}$.

 



\textbf{Solution 2}

Since $.48017$ is a bit less than $.5$, $(.48017)^3$ should be a bit less than $(.5)^3 = .125$.  Of the answer choices, $\boxed{\textbf{(B)}\ 0.110}$ is a bit less than $.125$ so it is the closest.

 


	\item (\textit{AMC 8 1990}) Which of these five numbers is the largest?


 $\text{(A)}\ 13579+\frac{1}{2468} \qquad \text{(B)}\ 13579-\frac{1}{2468} \qquad \text{(C)}\ 13579\times \frac{1}{2468}$
$\text{(D)}\ 13579\div \frac{1}{2468} \qquad \text{(E)}\ 13579.2468$

 \textbf{Answer: }D



\textbf{Solution}

Choices (A) and (E) are slightly more than $13579$.

 Choice (B) is slightly less than $13579$.

 Choice (C), which is equal to $\frac{13579}{2468}$, is much less than $13579$.

 Choice (D), which is equal to $13597 \times 2468$, is much larger than $13579$.

 So $\boxed{\textbf{(D)}\ 13579\div \frac{1}{2468}}$ is by far the largest of the answer choices.

 


	\item (\textit{AMC 8 1990}) When three different numbers from the set $\{ -3, -2, -1, 4, 5 \}$ are multiplied, the largest possible product is 


 $\text{(A)}\ 10 \qquad \text{(B)}\ 20 \qquad \text{(C)}\ 30 \qquad \text{(D)}\ 40 \qquad \text{(E)}\ 60$

 \textbf{Answer: }C



\textbf{Solution}

For the product of $3$ real numbers to be positive, either all $3$ numbers are positive, or $1$ number is positive and $2$ numbers are negative.

 Since there are only $2$ positive numbers, the only way to get a positive product is to select $2$ negative numbers and $1$ positive number.  The product is maximized by selecting the $2$ negative numbers with the largest absolute values and the positive number with the largest absolute value.  In this case the product would be $(-3)(-2)(5) = \boxed{\textbf{(C)}\ 30}$.

 


	\item (\textit{AMC 8 1990}) A dress originally priced at $80$ dollars was put on sale for $25\%$ off.  If $10\%$ tax was added to the sale price, then the total selling price (in dollars) of the dress was


 $\text{(A)}\ \text{45 dollars} \qquad \text{(B)}\ \text{52 dollars} \qquad \text{(C)}\ \text{54 dollars} \qquad \text{(D)}\ \text{66 dollars} \qquad \text{(E)}\ \text{68 dollars}$

 \textbf{Answer: }D



\textbf{Solution}

The sale price is $0.75 \cdot \$80 = \$60$, so the total price is $1.1 \cdot \$60 = \boxed{\textbf{(D)}\ \$66}$.

 


	\item (\textit{AMC 8 1990}) The grading scale shown is used at Jones Junior High.  The fifteen scores in Mr. Freeman's class were: 
 \[\begin{tabular}[t]{lllllllll} 89, & 72, & 54, & 97, & 77, & 92, & 85, & 74, & 75, \\ 63, & 84, & 78, & 71, & 80, & 90. & & & \\ \end{tabular}\]


 In Mr. Freeman's class, what percent of the students received a grade of C?


 
 \[\boxed{\begin{tabular}[t]{cc} A: & 93 - 100 \\ B: & 85 - 92 \\ C: & 75 - 84 \\ D: & 70 - 74 \\ F: & 0 - 69  \end{tabular}}\]
$\text{(A)}\ 20\% \qquad \text{(B)}\ 25\% \qquad \text{(C)}\ 30\% \qquad \text{(D)}\ 33\frac{1}{3}\% \qquad \text{(E)}\ 40\%$

 \textbf{Answer: }D



\textbf{Solution}

Of the $15$ students in Mr. Freeman's class, $5$ students (those with scores of $77$, $75$, $84$, $78$, and $80$) received a grade of C.  So $\frac{5}{15}=\boxed{\textbf{(D)}\ 33\frac{1}{3}\%}$ of his students received a grade of C.

 


	\item (\textit{AMC 8 1990}) On this monthly calendar, the date behind one of the letters is added to the date behind $\text{C}$.  If this sum equals the sum of the dates behind $\text{A}$ and $\text{B}$, then the letter is


 
                    \begin{figure}[H]
                        \centering
                        \includegraphics[width=0.5\textwidth]{Problem/AMC_8_AJHSME_1990_10_Problem_Figure_1.png}
                    \end{figure}
                    
$\text{(A)}\ \text{P} \qquad \text{(B)}\ \text{Q} \qquad \text{(C)}\ \text{R} \qquad \text{(D)}\ \text{S} \qquad \text{(E)}\ \text{T}$

 \textbf{Answer: }A



\textbf{Solution 1}

Since the date behind C is $1$ less than the date behind A, the date added to the date behind C should be $1$ more than the date behind B, which would be the date behind $\boxed{\textbf{(A)}\ \text{P}}$.

 



\textbf{Solution 2}

Let the date behind C be $x$, and $y$ be the date added to C. Then the date behind A is $x+1$, while the date behind B is $x + 13$.  So $(x+1)+(x+13)=x+y$, and therefore $y=x+14$, which is the date behind $\boxed{\textbf{(A)}\ \text{P}}$. 

 -motorfinn, edited by j314andrews

 


	\item (\textit{AMC 8 1990}) The numbers on the faces of this cube are consecutive whole numbers.  The sums of the two numbers on each of the three pairs of opposite faces are equal.  The sum of the six numbers on this cube is


 
                    \begin{figure}[H]
                        \centering
                        \includegraphics[width=0.5\textwidth]{Problem/AMC_8_AJHSME_1990_11_Problem_Figure_1.png}
                    \end{figure}
                    
$\text{(A)}\ 75 \qquad \text{(B)}\ 76 \qquad \text{(C)}\ 78 \qquad \text{(D)}\ 80 \qquad \text{(E)}\ 81$

 \textbf{Answer: }E



\textbf{Solution}

Since $11$ and $15$ are among the $6$ consecutive whole numbers, the numbers must be either $\{11,12,13,14,15,16\}$ or $\{10,11,12,13,14,15\}$.

 If the numbers were $\{10,11,12,13,14,15\}$, each pair of opposite faces would sum to $\frac{10+11+12+13+14+15}{3}=25$.  But $11$ would be opposite $14$, which is not the case on this cube.

 Therefore, the numbers are $\{11,12,13,14,15,16\}$, which works if $16$, $13$, and $12$ are placed opposite $11$, $14$, and $15$ respectively.  In this case, the sum is $11+12+13+14+15+16=\boxed{\textbf{(E)}\ 81}$.

 


	\item (\textit{AMC 8 1990}) There are twenty-four $4$-digit numbers that use each of the four digits $2$, $4$, $5$, and $7$ exactly once.  Listed in numerical order from smallest to largest, the number in the $17\text{th}$ position in the list is


 $\text{(A)}\ 4527 \qquad \text{(B)}\ 5724 \qquad \text{(C)}\ 5742 \qquad \text{(D)}\ 7245 \qquad \text{(E)}\ 7524$

 \textbf{Answer: }B



\textbf{Solution}

Each digit appears as the thousands digit in $3!=6$ numbers. Therefore, the $13\text{th}$ through $18\text{th}$ numbers start with $5$, and thus the $17\text{th}$ number on the list is the $5\text{th}$ number starting with $5$.  The numbers starting with $5$ from smallest to largest, are $5247,5274,5427,5472,5724,5742$, so $\boxed{\textbf{(B)}\ 5724}$ is $17\text{th}$ on the list.

 


	\item (\textit{AMC 8 1990}) One proposal for new postage rates for a letter was $30$ cents for the first ounce and $22$ cents for each additional ounce (or fraction of an ounce).  The postage for a letter weighing $4.5$ ounces was


 $\text{(A)}\ \text{96 cents} \qquad \text{(B)}\ \text{1.07 dollars} \qquad \text{(C)}\ \text{1.18 dollars} \qquad \text{(D)}\ \text{1.20 dollars} \qquad \text{(E)}\ \text{1.40 dollars}$

 \textbf{Answer: }C



\textbf{Solution}

There are $4.5-1=3.5$ additional ounces, which would round up to paying for $4$ additional ounces.  Therefore, the total cost is $\$0.30 + 4 \cdot \$ 0.22 = \boxed{\textbf{(B)}\ \$1.18}$.

 


	\item (\textit{AMC 8 1990}) A bag contains only blue balls and green balls.  There are $6$ blue balls.  If the probability of drawing a blue ball at random from this bag is $\frac{1}{4}$, then the number of green balls in the bag is


 $\text{(A)}\ 12 \qquad \text{(B)}\ 18 \qquad \text{(C)}\ 24 \qquad \text{(D)}\ 30 \qquad \text{(E)}\ 36$

 \textbf{Answer: }B



\textbf{Solution 1}

Since $\frac{1}{4}$ of the balls are blue, there are $4\times 6=24$ total balls. Thus $24-6=\boxed{\textbf{(B)}\ 18}$ balls are green.

 



\textbf{Solution 2}

Let $x$ be the number of green balls in the bag.  Then the probability of selecting a blue ball is $\frac{6}{6+x}=\frac{1}{4}$.  Solving this equation yields $x = \boxed{\textbf{(B)}\ 18}$.

 - goldenn, edited by j314andrews

 


	\item (\textit{AMC 8 1990}) The area of this figure is $100\text{ cm}^2$.  Its perimeter is 


 
                    \begin{figure}[H]
                        \centering
                        \includegraphics[width=0.5\textwidth]{Problem/AMC_8_AJHSME_1990_15_Problem_Figure_1.png}
                    \end{figure}
                    


 [figure consists of four identical squares]

 $\text{(A)}\ \text{20 cm} \qquad \text{(B)}\ \text{25 cm} \qquad \text{(C)}\ \text{30 cm} \qquad \text{(D)}\ \text{40 cm} \qquad \text{(E)}\ \text{50 cm}$

 \textbf{Answer: }E



\textbf{Solution}

Since the area of the whole figure is $100\text{ cm}^2$, each square has an area of $\frac{100}{4} = 25\text{ cm}^2$ and a side length of $\sqrt{25} = 5\text{ cm}$.  

 The boundary of this figure is made of $10$ segments, each of length $5\text{ cm}$, so it perimeter is $10 \cdot 5=\boxed{\textbf{(E)}\ 50\text{ cm}}$.

 


	\item (\textit{AMC 8 1990}) $1990-1980+1970-1960+\cdots -20+10 =$
$\text{(A)}\ -990 \qquad \text{(B)}\ -10 \qquad \text{(C)}\ 990 \qquad \text{(D)}\ 1000 \qquad \text{(E)}\ 1990$

 \textbf{Answer: }D



\textbf{Solution 1}

This sum has $\frac{1990}{10} = 199$ terms.  Grouping the sum as $(1990-1980)+(1970-1960)+...+(30-20)+10$ creates $\frac{198}{2} + 1 = 100$ groups which each sum to $10$, so the total sum is $100\cdot 10 = \boxed{\textbf{(D)}\ 1000}$.

 



\textbf{Solution 2}

As in Solution 1, this sum has $199$ terms.  Grouping the sum as $1990 + (-1980+1970) + (-1960+1950) + ... + (-20+10)$ creates $\frac{198}{2} = 99$ groups which each sum to $-10$, so the total sum is $1990 + (99)(-10) = \boxed{\textbf{(D)}\ 1000}$.

 



\textbf{Solution 3}

Grouping the terms as $(1990 + 10) - (1980 + 20) + ... + (1030 + 970) - (1020 + 980) + (1010 + 990) - 1000$ shows that the sum is $(2000 - 2000) + ... + (2000 - 2000) + 2000 - 10000 = 0 + ... + 0 + 1000 = \boxed{\textbf{(D)}\ 1000}$.

 


 


	\item (\textit{AMC 8 1990}) A straight concrete sidewalk is to be $3$ feet wide, $60$ feet long, and $3$ \textbf{inches} thick.  How many cubic yards of concrete must a contractor order for the sidewalk if concrete must be ordered in a whole number of cubic yards?


 $\text{(A)}\ 2 \qquad \text{(B)}\ 5 \qquad \text{(C)}\ 12 \qquad \text{(D)}\ 20 \qquad \text{(E)}\ \text{more than 20}$

 \textbf{Answer: }A



\textbf{Solution}

The sidewalk is $1$ yard wide, $20$ yards long, and $\frac{1}{12}$ yard thick, so its volume is $1\cdot 20\cdot\frac{1}{12} = 1\frac{2}{3}$ cubic yards.  So $\boxed{\textbf{(A)}\ 2}$ cubic yards must be ordered.

 


	\item (\textit{AMC 8 1990}) Each corner of a rectangular prism is cut off.  Two (of the eight) cuts are shown.  How many edges does the new figure have?


 
                    \begin{figure}[H]
                        \centering
                        \includegraphics[width=0.5\textwidth]{Problem/AMC_8_AJHSME_1990_18_Problem_Figure_1.png}
                    \end{figure}
                    
$\text{(A)}\ 24 \qquad \text{(B)}\ 30 \qquad \text{(C)}\ 36 \qquad \text{(D)}\ 42 \qquad \text{(E)}\ 48$

 \textbf{Answer: }C



\textbf{Solution}

The new figure retains the middle part of each of the $12$ edges of the cube and gains $3$ new edges from each cut.  Therefore, the new figure has $12 + 3 \cdot 8 = \boxed{\textbf{(C)}\ 36}$ edges.

 


	\item (\textit{AMC 8 1990}) There are $120$ seats in a row.  What is the fewest number of seats that must be occupied so the next person to be seated must sit next to someone?


 $\text{(A)}\ 30 \qquad \text{(B)}\ 40 \qquad \text{(C)}\ 41 \qquad \text{(D)}\ 60 \qquad \text{(E)}\ 119$

 \textbf{Answer: }B



\textbf{Solution 1}

A qualifying arrangement with $40$ occupied seats can be created by dividing the $120$-seat row into $40$ groups of $3$ adjacent seats and placing a person in the middle seat of each group.  

 It is impossible to find a qualifying arrangement with $39$ or fewer occupied seats.  Suppose the row of seats is initially empty, and each person who sits down must choose an empty seat that is not next to an occupied seat.  Then each person who sits down will eliminate at most $3$ available seats from people who sit down later.  Therefore, after $39$ people sit down, there must be at least $120 - 3 \cdot 39 = 3$ available seats, and thus a qualifying arrangement with $39$ or fewer occupied seats does not exist.

 Therefore, a qualifying arrangement has at least $\boxed{\textbf{(B)}\ 40}$ occupied seats.

 



\textbf{Solution 2}

Consider a row of chairs with $n$ occupied seats such that each empty seat is next to an occupied seat.  Then there are at most $2$ empty seats between any of the $n-1$ pairs of consecutive occupied seats, at most $1$ empty seat before the first occupied seat and at most $1$ empty seat after the last occupied seat.  Therefore, the row must have at most $n + 2(n-1) + 2 = 3n$ chairs in total.  

 Therefore, a row with $39$ or fewer occupied seats may have at most $39 \cdot 3 = 117$ seats, which is too few.  A row with $\boxed{\textbf{(B)}\ 40}$ occupied seats and $120$ total seats can be created by placing $1$ empty seat on each end and $2$ empty seats between each pair of consecutive occupied seats.

 


	\item (\textit{AMC 8 1990}) The annual incomes of $1,000$ families range from $8200$ dollars to $98,000$ dollars.  In error, the largest income was entered on the computer as $980,000$ dollars.  The difference between the mean of the incorrect data and the mean of the actual data is


 $\text{(A)}\ \text{882 dollars} \qquad \text{(B)}\ \text{980 dollars} \qquad \text{(C)}\ \text{1078 dollars} \qquad \text{(D)}\ \text{482,000 dollars} \qquad \text{(E)}\ \text{882,000 dollars}$

 \textbf{Answer: }A



\textbf{Solution 1}

The sum of the incorrectly entered incomes is $\$980000-\$98000=\$882000$ higher than the sum of the actual incomes.  Therefore, the mean of the incorrect incomes is $\frac{\$882000}{1000} = \boxed{\textbf{(A)}\ \$882}$ higher than the mean of the actual incomes.

 



\textbf{Solution 2}

Let $S$ be the sum of the correctly entered incomes.  The mean of the actual incomes is $\frac{S+98000}{1000}$, and for the incorrect data the mean is $\frac{S+980000}{1000}$.  The difference between these means is $\frac{S+980000}{1000}-\frac{S+98000}{1000} = \boxed{\textbf{(A)}\ \$882}$.

 


	\item (\textit{AMC 8 1990}) A list of $8$ numbers is formed by beginning with two given numbers.  Each new number in the list is the product of the two previous numbers.  Find the first number if the last three are shown:

 \[\text{\underline{\hspace{3 mm}?\hspace{3 mm}}\hspace{1 mm},\hspace{1 mm} \underline{\hspace{7 mm}}\hspace{1 mm},\hspace{1 mm} \underline{\hspace{7 mm}}\hspace{1 mm},\hspace{1 mm} \underline{\hspace{7 mm}}\hspace{1 mm},\hspace{1 mm} \underline{\hspace{7 mm}}\hspace{1 mm},\hspace{1 mm}\underline{\hspace{2 mm}16\hspace{2 mm}}\hspace{1 mm},\hspace{1 mm}\underline{\hspace{2 mm}64\hspace{2 mm}}\hspace{1 mm},\hspace{1 mm}\underline{\hspace{1 mm}1024\hspace{1 mm}}}\]
$\text{(A)}\ \frac{1}{64} \qquad \text{(B)}\ \frac{1}{4} \qquad \text{(C)}\ 1 \qquad \text{(D)}\ 2 \qquad \text{(E)}\ 4$

 \textbf{Answer: }B



\textbf{Solution}

Each missing term must be the quotient of the term two entries after it and the term immediately after it.  Therefore, the fifth term is $4$, the fourth term is $4$, the third is $1$, the second is $4$, and the first is $\boxed{\textbf{(B)}\ \frac{1}{4}}$.

 


	\item (\textit{AMC 8 1990}) Several students are seated at a large circular table.  They pass around a bag containing $100$ pieces of candy.  Each person receives the bag, takes one piece of candy and then passes the bag to the next person.  If Chris takes the first and last piece of candy, then the number of students at the table could be


 $\text{(A)}\ 10 \qquad \text{(B)}\ 11 \qquad \text{(C)}\ 19 \qquad \text{(D)}\ 20 \qquad \text{(E)}\ 25$

 \textbf{Answer: }B



\textbf{Solution}

Since the bag had $100$ candies at first, it must have been passed $100 - 1 = 99$ times.  Since the bag started and ended at the same student, the number of students at the table must be a factor of $99$.  Of the answer choices, only $\boxed{\textbf{(B)}\ 11}$ is a factor of $99$.

 


	\item (\textit{AMC 8 1990}) The graph relates the distance traveled [in miles] to the time elapsed [in hours] on a trip taken by an experimental airplane.  During which hour was the average speed of this airplane the largest?


 
                    \begin{figure}[H]
                        \centering
                        \includegraphics[width=0.5\textwidth]{Problem/AMC_8_AJHSME_1990_23_Problem_Figure_1.png}
                    \end{figure}
                    
$\text{(A)}\ \text{first (0-1)} \qquad \text{(B)}\ \text{second (1-2)} \qquad \text{(C)}\ \text{third (2-3)} \qquad \text{(D)}\ \text{ninth (8-9)} \qquad \text{(E)}\ \text{last (11-12)}$

 \textbf{Answer: }B



\textbf{Solution}

The time when the average speed is greatest is when the slope of the graph is steepest.  This is in the second hour $\rightarrow \boxed{\text{B}}$.

 


	\item (\textit{AMC 8 1990}) Three $\Delta$'s and a $\diamondsuit$ will balance nine $\bullet$'s.  One $\Delta$ will balance a $\diamondsuit$ and a $\bullet$.


 
                    \begin{figure}[H]
                        \centering
                        \includegraphics[width=0.5\textwidth]{Problem/AMC_8_AJHSME_1990_24_Problem_Figure_1.png}
                    \end{figure}
                    


 How many $\bullet$'s will balance the two $\diamondsuit$'s in this balance?


 
                    \begin{figure}[H]
                        \centering
                        \includegraphics[width=0.5\textwidth]{Problem/AMC_8_AJHSME_1990_24_Problem_Figure_2.png}
                    \end{figure}
                    
$\text{(A)}\ 1 \qquad \text{(B)}\ 2 \qquad \text{(C)}\ 3 \qquad \text{(D)}\ 4 \qquad \text{(E)}\ 5$

 \textbf{Answer: }C



\textbf{Solution 1}

Consider the scale in which $3$ $\triangle$'s and $1$ $\diamondsuit$ balance $9$ $\bullet$'s.

 After replacing each $\triangle$ on this scale with $1$ $\bullet$ and $1$ $\diamondsuit$, the scale will remain balanced.  This scale now has $4$ $\diamondsuit$'s and $3$ $\bullet$'s balancing $9$ $\bullet$'s.

 After removing $3$ $\bullet$'s from each side of the scale will remain balanced.  This scale now has $4$ $\diamondsuit$'s balancing $6$ $\bullet$'s.

 Finally, if each side of the scale is divided by $2$, the scale will remain balanced.  Therefore, $2$ $\diamondsuit$'s balances $\boxed{\textbf{(C)}\ 3}$ $\bullet$'s.

 


 



\textbf{Solution 2}

For simplicity, suppose $\Delta = a$, $\diamondsuit = b$ and $\bullet = c$.  Then,
 \[3a+b=9c\] 
 \[a=b+c\]and we want to know what $2b$ is in terms of $c$.  Substituting the second equation into the first, we have 
 \[4b=6c\Rightarrow 2b=3c\]

 Thus, we need $3$ $\bullet$'s $\rightarrow \boxed{\text{C}}$.

 


	\item (\textit{AMC 8 1990}) How many different patterns can be made by shading exactly two of the nine squares?  Patterns that can be matched by flips and/or turns are not considered different.  For example, the patterns shown below are \textbf{not} considered different.


 
                    \begin{figure}[H]
                        \centering
                        \includegraphics[width=0.5\textwidth]{Problem/AMC_8_AJHSME_1990_25_Problem_Figure_1.png}
                    \end{figure}
                    
$\text{(A)}\ 3 \qquad \text{(B)}\ 6 \qquad \text{(C)}\ 8 \qquad \text{(D)}\ 12 \qquad \text{(E)}\ 18$

 \textbf{Answer: }C



\textbf{Solution 1}

We break this into cases.

 \textit{Case 1: At least one square is a vertex:} \href{https://artofproblemsolving.com/wiki/index.php?title=WLOG}{WLOG}, suppose one of them is in the upper-left corner.  Then, consider the diagonal through that square.  The two squares on that diagonal could be the second square, or the second square is on one side of the diagonal.  

 The square is reflectionally symmetric about this diagonal, so we only consider the squares on one side, giving another three possibilities.  

 In this case, there are $2+3=5$ distinct squares.

 \textit{Case 2: At least one square is on an edge, but no square is on a vertex:} There are clearly two edge-edge combinations and one edge-center combination, so this case has $3$ squares.

 In total, there are $3+5=8$ distinct squares $\rightarrow \boxed{\text{C}}$




\textbf{Solution 2 (Burnside's Lemma)}

Consider each of the $8$ symmetries of the square.

 The identity transformation preserves all $\binom{9}{2} = 36$ arrangements.

 The $90^\circ$ rotations do not preserve any arrangements.

 The $180^\circ$ rotation preserves $4$ arrangements: those with either opposite corners or the centers of opposite edges shaded.

 The horizontal and vertical reflections each preserve $6$ arrangements: the $3$ in which the shaded squares are mirror images of each other and the $3$ in which both shaded squares are on the line of reflection.

 The diagonal reflections each preserve $6$ arrangements: the $3$ in which the shaded squares are mirror images of each other and the $3$ in which both shaded squares are on the line of reflection.

 Therefore, by Burnside's lemma, there are $\frac{1}{8}(36 + 4 + 2\cdot 6 + 2\cdot 6) = \boxed{\textbf{(C)}\ 8}$ different configurations if reflections and rotations are considered the same.

 -j314andrews

 


\end{enumerate}